% ============================================================
%  SOAL MATEMATIKA  --  SSU-ITB 2026
%  INTEN Set 1--3  (nomor 1--30)
%  Paket 3.1
% ============================================================

\documentclass[11pt]{article}

\usepackage[utf8]{inputenc}
\usepackage[T1]{fontenc}
\usepackage[a4paper,margin=2.5cm,top=2.8cm]{geometry}
\usepackage{amsmath,amssymb}
\usepackage{graphicx}
\usepackage{enumitem}
\usepackage{multicol}
\usepackage{xcolor}
\usepackage[most]{tcolorbox}
\usepackage{fancyhdr}
\usepackage{booktabs}
\usepackage{icomma}
\usepackage{array}

\pagestyle{fancy}\fancyhf{}
\lhead{\small SSU-ITB 2026}
\rhead{\small Paket 3.1}
\cfoot{\small\thepage}
\renewcommand{\headrulewidth}{0.4pt}

\newtcolorbox{soal}[1]{
  breakable,
  colback=white, colframe=black, colbacktitle=black,
  coltitle=white, fonttitle=\bfseries\sffamily,
  title=Nomor #1,
  boxrule=0.8pt, arc=0pt,
  left=3mm, right=3mm, top=2mm, bottom=2mm}

\newenvironment{pil}{%
  \begin{enumerate}[label=(\Alph*),leftmargin=*,itemsep=1pt,topsep=3pt]}%
  {\end{enumerate}}
\newenvironment{pilB}{%
  \par\smallskip\begin{multicols}{2}
  \begin{enumerate}[label=(\Alph*),leftmargin=*,itemsep=1pt,topsep=3pt]}%
  {\end{enumerate}\end{multicols}}

\newcommand{\gambar}[2][0.52\textwidth]{%
  \begin{center}\includegraphics[width=#1]{#2}\end{center}}

\linespread{1.05}

% ===================================================================
\begin{document}
\thispagestyle{empty}

\begin{center}
  \fbox{\parbox{0.45\textwidth}{\centering\bfseries\large LEMBAR SOAL}}
\end{center}
\vspace{0.5cm}

\begin{center}
  {\LARGE\bfseries MATEMATIKA}
\end{center}
\vspace{0.5cm}

\begin{center}
  \begin{tabular}{@{\hspace{2cm}}ll@{\hspace{0.4cm}:@{\hspace{0.4cm}}}l}
    Jenjang      & & SMA / Sederajat               \\[0.3em]
    Hari/Tanggal & & \underline{\hspace{5cm}}      \\[0.3em]
    Waktu        & & \underline{\hspace{5cm}}      \\[0.3em]
    Paket Soal   & & 3.1                           \\
  \end{tabular}
\end{center}

\vspace{0.6cm}
\noindent\rule{\linewidth}{0.6pt}
\vspace{0.3cm}

\noindent\textbf{\underline{PETUNJUK UMUM}}
\begin{enumerate}[leftmargin=*,itemsep=4pt,topsep=4pt]
  \item Anda \textbf{tidak} diperbolehkan menggunakan sumber-sumber eksternal,
        termasuk internet, \textit{large language model} seperti ChatGPT, Claude,
        LLaMA, Gemini, Meta AI, dan sebagainya. Anda sangat dianjurkan untuk
        mencoba mengerjakan sendiri terlebih dahulu karena evaluasi ini dibuat
        untuk \textbf{mengukur} tingkat kepahaman; nilai $<70$ mengindikasikan
        \textbf{tidak lulus}, dan jika sudah melewati \textit{deadline} maka
        dianggap tidak mengerjakan yang berarti \textbf{tidak lulus}.

  \item Terdapat 2 tahap pengerjaan, yaitu tahap \textbf{Try Out} dan tahap
        \textbf{Review}. Try Out bersifat \textit{open book}; Anda diperbolehkan
        membuka buku apapun selama bersifat \textit{offline} dengan rentang waktu
        2 jam pada tanggal \underline{\hspace{3cm}}. Setelah itu memasuki
        tahap Review dengan \textit{schedule} rentang tanggal
        \underline{\hspace{3cm}}--\underline{\hspace{3cm}}, di mana Anda
        diharapkan mengerjakan ulang dengan disertakan penjelasan/pembelaan
        mengapa revisi jawaban adalah pilihan tersebut, dan tetap mencantumkan
        penjelasan pada soal yang walaupun pada penilaian sudah benar.

  \item Anda dianjurkan untuk mencantumkan caranya walaupun tidak termasuk
        penilaian, mengingat sifatnya adalah sebagai bahan pembelajaran.
        Mohon bertanggung jawab atas pekerjaan Anda sendiri.

  \item Bacalah setiap soal dengan teliti. Terdapat tiga jenis soal:
        \begin{enumerate}[label=(\arabic*),leftmargin=2em,topsep=2pt,itemsep=2pt]
          \item \textbf{Pilihan Ganda Biasa} --- 5 opsi (A--E), pilih
                1 jawaban paling tepat;
          \item \textbf{Pilihan Ganda Majemuk Kompleks} --- tabel
                \textbf{Benar}/\textbf{Salah}, tentukan kebenaran tiap
                pernyataan secara mandiri; dan
          \item \textbf{Isian Singkat} --- tulis nilai/ekspresi pada
                kolom yang tersedia.
        \end{enumerate}

  \item Soal berjumlah \textbf{30 butir} soal dengan bobot asli per soal adalah
        $\mathbf{\tfrac{10}{3}}$ untuk mendapatkan nilai sempurna sebesar \textbf{100}.

  \item Silakan bertanya apabila terdapat soal yang ambigu, rancu, atau kurang
        spesifik selama itu tidak menjurus kepada jawaban soal tersebut.

  \item Isilah detail informasi berikut.\\[0.3em]
        \textbf{Nama}\hspace{0.45cm}:\enspace\underline{\hspace{8cm}}
\end{enumerate}

\vspace{0.6cm}
\noindent\textit{Dengan menandatangani kolom di bawah ini, saya menyatakan telah membaca,
memahami, dan menyetujui seluruh peraturan yang tercantum di atas.}

\vspace{0.6cm}
\noindent\fbox{\parbox{5cm}{\vspace{2.2cm}\centering\small Tanda Tangan Peserta\vspace{0.3cm}}}

\clearpage

% ===================================================================
%  SET 1  (nomor 1--10)
% ===================================================================

%---------- 1: PG Biasa ----------
\begin{soal}{1}
Jika $9^{1/6}\cdot x=3$, maka nilai $x$ adalah \ldots
\begin{pilB}
  \item $3^{4/6}$ \item $3^{4/36}$ \item $3^{6/2}$ \item $3^{6}$ \item $3^{1/3}$
\end{pilB}
\end{soal}

%---------- 2: PG Biasa ----------
\begin{soal}{2}
Jika titik $A(4,15)$ terletak pada grafik fungsi $g(x)=\dfrac{2x+7}{x^2-b}$,
maka nilai $b$ adalah \ldots
\begin{pilB}
  \item $16$ \item $14$ \item $13$ \item $15$ \item $12$
\end{pilB}
\end{soal}

%---------- 3: PG Majemuk Kompleks ----------
\begin{soal}{3}
Diketahui $x_0,x_1,x_2,\ldots,x_{100}$ adalah barisan aritmetika dengan
$x_0=3$ dan $x_{100}=9$. Tentukan kebenaran setiap pernyataan berikut!

\par\smallskip
{\renewcommand{\arraystretch}{1.4}%
\begin{tabular}{@{}p{0.76\linewidth}cc@{}}
\hline
\textbf{Pernyataan} & \textbf{Benar} & \textbf{Salah}\\
\hline
(1)~Beda barisan aritmetika tersebut adalah $\dfrac{3}{50}$.
  & $\bigcirc$ & $\bigcirc$\\
(2)~Nilai $x_{50} = 6$.
  & $\bigcirc$ & $\bigcirc$\\
(3)~Nilai $x_{25} = 4{,}5$.
  & $\bigcirc$ & $\bigcirc$\\
(4)~Nilai $x_{75} = 8$.
  & $\bigcirc$ & $\bigcirc$\\
\hline
\end{tabular}}
\end{soal}

%---------- 4: PG Majemuk Kompleks ----------
\begin{soal}{4}
Diberikan fungsi $f(x)=\dfrac{5x^2-9x+3}{(x-1)^6}$.
Tentukan kebenaran setiap pernyataan berikut!

\par\smallskip
{\renewcommand{\arraystretch}{1.4}%
\begin{tabular}{@{}p{0.76\linewidth}cc@{}}
\hline
\textbf{Pernyataan} & \textbf{Benar} & \textbf{Salah}\\
\hline
(1)~$f'(x)$ terdefinisi untuk semua $x\neq1$.
  & $\bigcirc$ & $\bigcirc$\\
(2)~$f'(2) = -19$.
  & $\bigcirc$ & $\bigcirc$\\
(3)~$f'(2) > 0$.
  & $\bigcirc$ & $\bigcirc$\\
(4)~Jika $g(x)=(x-1)^6 f(x)$, maka $g(2)=5$.
  & $\bigcirc$ & $\bigcirc$\\
\hline
\end{tabular}}
\end{soal}

%---------- 5: PG Biasa ----------
\begin{soal}{5}
Jika $x^3-2x+5$ dibagi oleh $1-x$, maka sisanya adalah \ldots
\begin{pilB}
  \item $4$ \item $5$ \item $6$ \item $7$ \item $8$
\end{pilB}
\end{soal}

%---------- 6: PG Biasa ----------
\begin{soal}{6}
Diketahui $f(x)=x^3+(a-1)x^2-2$. Jika $f(-2)=26$, maka nilai $a$ adalah \ldots
\begin{pilB}
  \item $5$ \item $10$ \item $15$ \item $20$ \item $25$
\end{pilB}
\end{soal}

%---------- 7: PG Majemuk Kompleks ----------
\begin{soal}{7}
Tiga buah dadu enam sisi dilempar bersamaan.
Tentukan kebenaran setiap pernyataan berikut!

\par\smallskip
{\renewcommand{\arraystretch}{1.4}%
\begin{tabular}{@{}p{0.76\linewidth}cc@{}}
\hline
\textbf{Pernyataan} & \textbf{Benar} & \textbf{Salah}\\
\hline
(1)~Ruang sampel pelemparan tiga dadu berjumlah $216$.
  & $\bigcirc$ & $\bigcirc$\\
(2)~Banyaknya cara agar jumlah tiga mata dadu bernilai $7$ adalah $15$.
  & $\bigcirc$ & $\bigcirc$\\
(3)~Peluang jumlah mata dadu ketiga bernilai $7$ adalah $\dfrac{15}{216}$.
  & $\bigcirc$ & $\bigcirc$\\
(4)~Peluang jumlah mata dadu ketiga bernilai $7$ adalah $\dfrac{1}{12}$.
  & $\bigcirc$ & $\bigcirc$\\
\hline
\end{tabular}}
\end{soal}

%---------- 8: PG Biasa ----------
\begin{soal}{8}
Diberikan trapesium $ABCD$ seperti pada gambar. Jika panjang $AB=11$,
$CD=7$, $BC=5$, dan luas trapesium adalah 27, maka keliling trapesium
adalah \ldots
\gambar[0.38\textwidth]{images/3/3-2.png}
\begin{pilB}
  \item $21$ \item $24$ \item $26$ \item $27$ \item $30$
\end{pilB}
\end{soal}

%---------- 9: PG Biasa ----------
\begin{soal}{9}
Diberikan persamaan lingkaran $(x+5)^2+(y-3)^2=9$. Titik pusat dan
jari-jari lingkaran tersebut adalah \ldots
\begin{pil}
  \item Titik pusat $(5,3)$ dan jari-jari $9$
  \item Titik pusat $(-5,3)$ dan jari-jari $3$
  \item Titik pusat $(5,3)$ dan jari-jari $3$
  \item Titik pusat $(5,-3)$ dan jari-jari $9$
  \item Titik pusat $(-5,-3)$ dan jari-jari $3$
\end{pil}
\end{soal}

%---------- 10: PG Majemuk Kompleks ----------
\begin{soal}{10}
Misalkan $\dfrac{A}{2-\sqrt{5}}+\dfrac{B}{2+\sqrt{5}}=1+2\sqrt{5}$
dengan $A$ dan $B$ bilangan real.
Tentukan kebenaran setiap pernyataan berikut!

\par\smallskip
{\renewcommand{\arraystretch}{1.4}%
\begin{tabular}{@{}p{0.76\linewidth}cc@{}}
\hline
\textbf{Pernyataan} & \textbf{Benar} & \textbf{Salah}\\
\hline
(1)~$A + B = -\dfrac{1}{2}$.
  & $\bigcirc$ & $\bigcirc$\\
(2)~$A = -\dfrac{5}{4}$.
  & $\bigcirc$ & $\bigcirc$\\
(3)~$B = -\dfrac{3}{4}$.
  & $\bigcirc$ & $\bigcirc$\\
(4)~$A < B$.
  & $\bigcirc$ & $\bigcirc$\\
\hline
\end{tabular}}
\end{soal}

% ===================================================================
%  SET 2  (nomor 11--20)
% ===================================================================

%---------- 11: PG Biasa ----------
\begin{soal}{11}
Jika sistem persamaan $\begin{cases}px+qy=8\\ 3x-qy=38\end{cases}$
memiliki penyelesaian $(x,y)=(2,4)$, maka nilai $p$ adalah \ldots
\begin{pilB}
  \item $40$ \item $22{,}5$ \item $21{,}5$ \item $20$ \item $8$
\end{pilB}
\end{soal}

%---------- 12: PG Biasa ----------
\begin{soal}{12}
Jika $a=\sqrt{7}+3$ dan $b=\sqrt{7}-3$, maka $a^3+b^3$ adalah \ldots
\begin{pilB}
  \item $60\sqrt{7}$ \item $62\sqrt{7}$ \item $64\sqrt{7}$
  \item $66\sqrt{7}$ \item $68\sqrt{7}$
\end{pilB}
\end{soal}

%---------- 13: PG Biasa ----------
\begin{soal}{13}
Jika $a<x<b$ adalah solusi dari $\dfrac{x^2+x+3}{x^2-x-2}<0$,
maka nilai $b-2a$ adalah \ldots
\begin{pilB}
  \item $1$ \item $2$ \item $3$ \item $4$ \item $5$
\end{pilB}
\end{soal}

%---------- 14: PG Majemuk Kompleks ----------
\begin{soal}{14}
Diberikan $A=\begin{pmatrix} a & b & c\\ 1 & -1 & 1\end{pmatrix}$
dan $B=\begin{pmatrix} 1 & 2\\ 1 & -1\\ 0 & 1\end{pmatrix}$.
Tentukan kebenaran setiap pernyataan berikut!

\par\smallskip
{\renewcommand{\arraystretch}{1.4}%
\begin{tabular}{@{}p{0.76\linewidth}cc@{}}
\hline
\textbf{Pernyataan} & \textbf{Benar} & \textbf{Salah}\\
\hline
(1)~Matriks $AB$ berukuran $2\times2$.
  & $\bigcirc$ & $\bigcirc$\\
(2)~$\det(AB) = 4(a+b)$.
  & $\bigcirc$ & $\bigcirc$\\
(3)~Jika $\det(AB)=4$, maka $a+b=4$.
  & $\bigcirc$ & $\bigcirc$\\
(4)~Jika $\det(AB)=4$, maka $a+b=1$.
  & $\bigcirc$ & $\bigcirc$\\
\hline
\end{tabular}}
\end{soal}

%---------- 15: PG Biasa ----------
\begin{soal}{15}
Agar persamaan kuadrat $ax^2+bx+b=0$ mempunyai dua akar real berbeda
untuk $b<0$ atau $b>8$, maka nilai $a$ adalah \ldots
\begin{pilB}
  \item $1$ \item $2$ \item $-2$ \item $3$ \item $-3$
\end{pilB}
\end{soal}

%---------- 16: PG Majemuk Kompleks ----------
\begin{soal}{16}
Diketahui $f(x)=3x^2+9x+12$. Jika
$\dfrac{f(x+h)-f(x)}{h}=Ax+B+Ch$,
tentukan kebenaran setiap pernyataan berikut!

\par\smallskip
{\renewcommand{\arraystretch}{1.4}%
\begin{tabular}{@{}p{0.76\linewidth}cc@{}}
\hline
\textbf{Pernyataan} & \textbf{Benar} & \textbf{Salah}\\
\hline
(1)~$A = 6$.
  & $\bigcirc$ & $\bigcirc$\\
(2)~$B = 3$.
  & $\bigcirc$ & $\bigcirc$\\
(3)~$C = 3$.
  & $\bigcirc$ & $\bigcirc$\\
(4)~$100A+10B+C = 693$.
  & $\bigcirc$ & $\bigcirc$\\
\hline
\end{tabular}}
\end{soal}

%---------- 17: PG Biasa ----------
\begin{soal}{17}
Misalkan $f(x)=\begin{cases}x^2-x+2, & x\geq1\\ x+1, & x<1\end{cases}$.
Maka nilai $\lim_{x\to1}f(x)$ adalah \ldots
\begin{pilB}
  \item $-2$ \item $-1$ \item Tidak ada \item $1$ \item $2$
\end{pilB}
\end{soal}

%---------- 18: PG Biasa ----------
\begin{soal}{18}
Jika $f(3)=6$, $f'(3)=4$, $g(3)=10$, dan $g'(3)=-6$, maka
$\left(\dfrac{f}{g}\right)'(3)$ adalah \ldots
\begin{pilB}
  \item $0{,}76$ \item $0{,}82$ \item $0{,}91$ \item $0{,}95$ \item $1$
\end{pilB}
\end{soal}

%---------- 19: PG Majemuk Kompleks ----------
\begin{soal}{19}
Diketahui $\displaystyle\int_1^9 f(x)\,dx=5$.
Tentukan kebenaran setiap pernyataan berikut!

\par\smallskip
{\renewcommand{\arraystretch}{1.4}%
\begin{tabular}{@{}p{0.76\linewidth}cc@{}}
\hline
\textbf{Pernyataan} & \textbf{Benar} & \textbf{Salah}\\
\hline
(1)~Substitusi $u=x^2$ mengubah batas integral $\displaystyle\int_1^3$
    menjadi $\displaystyle\int_1^9$.
  & $\bigcirc$ & $\bigcirc$\\
(2)~$\displaystyle\int_1^3 8x\,f(x^2)\,dx = 4\int_1^9 f(u)\,du$.
  & $\bigcirc$ & $\bigcirc$\\
(3)~Nilai $\displaystyle\int_1^3 8x\,f(x^2)\,dx = 5$.
  & $\bigcirc$ & $\bigcirc$\\
(4)~Nilai $\displaystyle\int_1^3 8x\,f(x^2)\,dx = 20$.
  & $\bigcirc$ & $\bigcirc$\\
\hline
\end{tabular}}
\end{soal}

%---------- 20: PG Majemuk Kompleks ----------
\begin{soal}{20}
Perhatikan pertidaksamaan $|x-7|\geq3$.
Tentukan kebenaran setiap pernyataan berikut!

\par\smallskip
{\renewcommand{\arraystretch}{1.4}%
\begin{tabular}{@{}p{0.76\linewidth}cc@{}}
\hline
\textbf{Pernyataan} & \textbf{Benar} & \textbf{Salah}\\
\hline
(1)~Bilangan bulat yang \textbf{tidak} memenuhi $|x-7|\geq3$
    adalah bilangan bulat yang memenuhi $|x-7|<3$.
  & $\bigcirc$ & $\bigcirc$\\
(2)~Kondisi $|x-7|<3$ setara dengan $4<x<10$.
  & $\bigcirc$ & $\bigcirc$\\
(3)~Banyaknya bilangan bulat yang tidak memenuhi pertidaksamaan tersebut
    adalah $4$.
  & $\bigcirc$ & $\bigcirc$\\
(4)~Banyaknya bilangan bulat yang tidak memenuhi pertidaksamaan tersebut
    adalah $5$.
  & $\bigcirc$ & $\bigcirc$\\
\hline
\end{tabular}}
\end{soal}

% ===================================================================
%  SET 3  (nomor 21--30)
% ===================================================================

%---------- 21: PG Biasa ----------
\begin{soal}{21}
Jika $9^m=4$, maka $4\cdot3^{m+1}-27^m$ adalah \ldots
\begin{pilB}
  \item $8$ \item $12$ \item $16$ \item $18$ \item $24$
\end{pilB}
\end{soal}

%---------- 22: PG Biasa ----------
\begin{soal}{22}
Nilai $a$ yang memenuhi
$\dfrac{1}{\log_{10}a}+\dfrac{1}{\log_{\sqrt{10}}a}+\dfrac{1}{\log_{\sqrt{\sqrt{10}}}a}+\cdots=200$
adalah \ldots
\begin{pilB}
  \item $\dfrac{1}{100}$ \item $\dfrac{1}{10}$ \item $\sqrt{10}$
  \item $10^{1/100}$ \item $10^{1/10}$
\end{pilB}
\end{soal}

%---------- 23: PG Majemuk Kompleks ----------
\begin{soal}{23}
Diberikan persamaan $\sin x=\dfrac{a-1{,}5}{2-0{,}5a}$.
Tentukan kebenaran setiap pernyataan berikut!

\par\smallskip
{\renewcommand{\arraystretch}{1.4}%
\begin{tabular}{@{}p{0.76\linewidth}cc@{}}
\hline
\textbf{Pernyataan} & \textbf{Benar} & \textbf{Salah}\\
\hline
(1)~Ruas kanan dapat ditulis sebagai $\dfrac{2a-3}{4-a}$.
  & $\bigcirc$ & $\bigcirc$\\
(2)~Agar persamaan mempunyai solusi, diperlukan
    $-1\leq\dfrac{2a-3}{4-a}\leq1$.
  & $\bigcirc$ & $\bigcirc$\\
(3)~Banyaknya bilangan bulat $a$ sehingga persamaan mempunyai
    penyelesaian adalah $3$.
  & $\bigcirc$ & $\bigcirc$\\
(4)~Banyaknya bilangan bulat $a$ sehingga persamaan mempunyai
    penyelesaian adalah $4$.
  & $\bigcirc$ & $\bigcirc$\\
\hline
\end{tabular}}
\end{soal}

%---------- 24: PG Majemuk Kompleks ----------
\begin{soal}{24}
Perhatikan $\displaystyle\sum_{k=1}^{100}2k+\sum_{k=1}^{100}(3k+2)$.
Tentukan kebenaran setiap pernyataan berikut!

\par\smallskip
{\renewcommand{\arraystretch}{1.4}%
\begin{tabular}{@{}p{0.76\linewidth}cc@{}}
\hline
\textbf{Pernyataan} & \textbf{Benar} & \textbf{Salah}\\
\hline
(1)~$\displaystyle\sum_{k=1}^{100} k = 4950$.
  & $\bigcirc$ & $\bigcirc$\\
(2)~$\displaystyle\sum_{k=1}^{100} 2k = 10100$.
  & $\bigcirc$ & $\bigcirc$\\
(3)~$\displaystyle\sum_{k=1}^{100} (3k+2) = 15350$.
  & $\bigcirc$ & $\bigcirc$\\
(4)~Total nilai $\displaystyle\sum_{k=1}^{100}2k+\sum_{k=1}^{100}(3k+2) = 25450$.
  & $\bigcirc$ & $\bigcirc$\\
\hline
\end{tabular}}
\end{soal}

%---------- 25: PG Majemuk Kompleks ----------
\begin{soal}{25}
Perhatikan pertidaksamaan $(x+1)(x^2+2x-7)\geq x^2-1$.
Tentukan kebenaran setiap pernyataan berikut!

\par\smallskip
{\renewcommand{\arraystretch}{1.4}%
\begin{tabular}{@{}p{0.76\linewidth}cc@{}}
\hline
\textbf{Pernyataan} & \textbf{Benar} & \textbf{Salah}\\
\hline
(1)~Pertidaksamaan dapat difaktorkan menjadi $(x-2)(x+1)(x+3)\geq0$.
  & $\bigcirc$ & $\bigcirc$\\
(2)~$x=-3$ memenuhi pertidaksamaan.
  & $\bigcirc$ & $\bigcirc$\\
(3)~Banyaknya bilangan bulat negatif yang memenuhi adalah $3$.
  & $\bigcirc$ & $\bigcirc$\\
(4)~$x=-4$ memenuhi pertidaksamaan.
  & $\bigcirc$ & $\bigcirc$\\
\hline
\end{tabular}}
\end{soal}

%---------- 26: PG Biasa ----------
\begin{soal}{26}
Hasil penjualan buku selama 10 hari terakhir bulan Juni adalah
$8,7,8,6,8,8,9,8,10,8$. Jika rata-rata penjualan pada bulan Juni
adalah 10, maka rata-rata penjualan pada 20 hari pertama adalah \ldots
\begin{pilB}
  \item $8$ \item $9$ \item $10$ \item $11$ \item $13$
\end{pilB}
\end{soal}

%---------- 27: PG Majemuk Kompleks ----------
\begin{soal}{27}
Diketahui $x^7-10x^6+24x^5=0$ memiliki tiga akar real berbeda
$r_1<r_2<r_3$.
Tentukan kebenaran setiap pernyataan berikut!

\par\smallskip
{\renewcommand{\arraystretch}{1.4}%
\begin{tabular}{@{}p{0.76\linewidth}cc@{}}
\hline
\textbf{Pernyataan} & \textbf{Benar} & \textbf{Salah}\\
\hline
(1)~Persamaan dapat difaktorkan menjadi $x^5(x-4)(x-6)=0$.
  & $\bigcirc$ & $\bigcirc$\\
(2)~Terdapat tiga akar real berbeda, yaitu $0,4,6$.
  & $\bigcirc$ & $\bigcirc$\\
(3)~$r_1+r_3 = 6$.
  & $\bigcirc$ & $\bigcirc$\\
(4)~$r_1\cdot r_3 = 24$.
  & $\bigcirc$ & $\bigcirc$\\
\hline
\end{tabular}}
\end{soal}

%---------- 28: PG Biasa ----------
\begin{soal}{28}
Jika persamaan garis singgung kurva $y=3x^2+C$ di suatu titik
adalah $y=24x$, maka nilai $C$ adalah \ldots
\begin{pilB}
  \item $50$ \item $46$ \item $47$ \item $48$ \item $49$
\end{pilB}
\end{soal}

%---------- 29: PG Majemuk Kompleks ----------
\begin{soal}{29}
Perhatikan transformasi parabola $y=x^2$.
Tentukan kebenaran setiap pernyataan berikut!

\par\smallskip
{\renewcommand{\arraystretch}{1.4}%
\begin{tabular}{@{}p{0.76\linewidth}cc@{}}
\hline
\textbf{Pernyataan} & \textbf{Benar} & \textbf{Salah}\\
\hline
(1)~Transformasi $y=(x-h)^2$ merupakan pergeseran grafik $y=x^2$
    sejauh $h$ satuan ke kanan (untuk $h>0$).
  & $\bigcirc$ & $\bigcirc$\\
(2)~Parabola $y=(x-1)^2$ diperoleh dengan menggeser $y=x^2$
    ke kanan 1 satuan.
  & $\bigcirc$ & $\bigcirc$\\
(3)~Parabola $y=(x-1)^2$ diperoleh dengan menggeser $y=x^2$
    ke kiri 1 satuan.
  & $\bigcirc$ & $\bigcirc$\\
(4)~Titik minimum parabola $y=(x-1)^2$ berada di $(1,0)$.
  & $\bigcirc$ & $\bigcirc$\\
\hline
\end{tabular}}
\end{soal}

%---------- 30: Isian ----------
\begin{soal}{30}
Bilangan bulat yang nilainya paling dekat dengan $5\sqrt{2}$ adalah \ldots\\[0.4em]
Jawaban:\enspace\underline{\hspace{4cm}}
\end{soal}

\end{document}
